% =====================================================================
% forge :: skeleton.tex  (v0.3)
% Vertical dependency-flow map as a READABLE LOGIC STORY. One page.
% Read down; each step follows from the ones above. Each NODE is a small
% two-line brick:
%   line 1:  <disc>  <FULL object name, bold via \ref>  --  <plain-English role>
%   line 2:  (indented, smaller) <one concrete detail>  [optional (needs <Name>)]
% Between steps: a centered down-arrow carrying the logic word ("so"/"then"/
% "still open"). Dependencies are implied by vertical position; an explicit
% "(needs <Full Name>)" is added ONLY for a cross-link that is not the node
% directly above. Numbers stay synced via \ref to the labels in prototype.tex.
% % =====================================================================
% forge :: skeleton.tex  (v0.3)
% Vertical dependency-flow map as a READABLE LOGIC STORY. One page.
% Read down; each step follows from the ones above. Each NODE is a small
% two-line brick:
%   line 1:  <disc>  <FULL object name, bold via \ref>  --  <plain-English role>
%   line 2:  (indented, smaller) <one concrete detail>  [optional (needs <Name>)]
% Between steps: a centered down-arrow carrying the logic word ("so"/"then"/
% "still open"). Dependencies are implied by vertical position; an explicit
% "(needs <Full Name>)" is added ONLY for a cross-link that is not the node
% directly above. Numbers stay synced via \ref to the labels in prototype.tex.
% % =====================================================================
% forge :: skeleton.tex  (v0.3)
% Vertical dependency-flow map as a READABLE LOGIC STORY. One page.
% Read down; each step follows from the ones above. Each NODE is a small
% two-line brick:
%   line 1:  <disc>  <FULL object name, bold via \ref>  --  <plain-English role>
%   line 2:  (indented, smaller) <one concrete detail>  [optional (needs <Name>)]
% Between steps: a centered down-arrow carrying the logic word ("so"/"then"/
% "still open"). Dependencies are implied by vertical position; an explicit
% "(needs <Full Name>)" is added ONLY for a cross-link that is not the node
% directly above. Numbers stay synced via \ref to the labels in prototype.tex.
% % =====================================================================
% forge :: skeleton.tex  (v0.3)
% Vertical dependency-flow map as a READABLE LOGIC STORY. One page.
% Read down; each step follows from the ones above. Each NODE is a small
% two-line brick:
%   line 1:  <disc>  <FULL object name, bold via \ref>  --  <plain-English role>
%   line 2:  (indented, smaller) <one concrete detail>  [optional (needs <Name>)]
% Between steps: a centered down-arrow carrying the logic word ("so"/"then"/
% "still open"). Dependencies are implied by vertical position; an explicit
% "(needs <Full Name>)" is added ONLY for a cross-link that is not the node
% directly above. Numbers stay synced via \ref to the labels in prototype.tex.
% \input{skeleton} from prototype.tex (after the appendix).
% =====================================================================
\clearpage
\section*{Dependency-flow skeleton}
\addcontentsline{toc}{section}{Dependency-flow skeleton}

% A NODE brick: disc, full name (bold), plain-English role; then an indented
% smaller concrete-detail line. #1 disc | #2 \ref name | #3 role | #4 detail.
\newcommand{\sknode}[4]{%
  \noindent #1\, \textbf{#2} \;---\; #3\par
  \nobreak\smallskip
  {\footnotesize\color{gray!75}\hspace*{1.6em}#4\par}}

% The logic arrow between steps, carrying the connective word.
\newcommand{\skstep}[1]{%
  \par\medskip\centerline{$\downarrow$\ \footnotesize\itshape\color{gray!70}#1}\par\medskip}

% A faint right-aligned tier tag (optional, must not dominate).
\newcommand{\sktag}[1]{\hfill{\scriptsize\color{gray!45}#1}}

\begin{center}
\begin{tcolorbox}[enhanced, breakable, width=0.92\linewidth,
                  colback=gray!3, colframe=gray!45, boxrule=0.5pt,
                  arc=3pt, left=12pt, right=12pt, top=8pt, bottom=8pt,
                  title={Argument flow \quad
                         \disc{v}\,V \;\disc{p}\,P \;\disc{h}\,H \;\disc{c}\,C}]

\sknode{\disc{v}}{Assumption~\ref{ass:cost}}{entry costs a little, and ties split evenly}{%
  $0<\kappa<1/n$ per front opened --- worth opening one, never worth opening all\sktag{setup}}

\skstep{so}

\sknode{\disc{p}}{Proposition~\ref{prop:exist}}{the contest has a mixed-strategy equilibrium}{%
  for any opened fronts, players randomize how hard to fight on each\sktag{engine}}

\skstep{and}

\sknode{\disc{p}}{Lemma~\ref{lem:concentrate}}{each side opens only a few fronts}{%
  at most $\lceil 1/(2\kappa)\rceil$ of them --- a bounded number, not all $n$\sktag{engine}}

\skstep{so}

\sknode{\disc{h}}{Corollary~\ref{cor:sparse}}{those few fronts don't grow with $n$}{%
  the support stays size $O(1/\kappa)$, independent of $n$\sktag{hinge}}

\skstep{so}

\sknode{\disc{c}}{Theorem~\ref{thm:char}}{exact closed-form for how they play}{%
  a symmetric mixed equilibrium: open $m^\star=\lceil 1/(2\kappa)\rceil$ fronts,
  spend uniformly on $[0,2/m^\star]$ each
  \;\textnormal{(needs Proposition~\ref{prop:exist})}\sktag{main result}}

\skstep{still open}

\sknode{\disc{c}}{Uniqueness?}{is this the only solution?}{%
  rule out asymmetric equilibria on the same support --- would turn
  Theorem~\ref{thm:char} green\sktag{open frontier}}

\medskip
\hrule
\smallskip
{\footnotesize\itshape
Read top-to-bottom: each step follows from the ones above; the word on each
$\downarrow$ is the logic. ``(needs $X$)'' flags a step that leans on something
not directly above it. Disc colour = current verification status.}

\end{tcolorbox}
\end{center}
 from prototype.tex (after the appendix).
% =====================================================================
\clearpage
\section*{Dependency-flow skeleton}
\addcontentsline{toc}{section}{Dependency-flow skeleton}

% A NODE brick: disc, full name (bold), plain-English role; then an indented
% smaller concrete-detail line. #1 disc | #2 \ref name | #3 role | #4 detail.
\newcommand{\sknode}[4]{%
  \noindent #1\, \textbf{#2} \;---\; #3\par
  \nobreak\smallskip
  {\footnotesize\color{gray!75}\hspace*{1.6em}#4\par}}

% The logic arrow between steps, carrying the connective word.
\newcommand{\skstep}[1]{%
  \par\medskip\centerline{$\downarrow$\ \footnotesize\itshape\color{gray!70}#1}\par\medskip}

% A faint right-aligned tier tag (optional, must not dominate).
\newcommand{\sktag}[1]{\hfill{\scriptsize\color{gray!45}#1}}

\begin{center}
\begin{tcolorbox}[enhanced, breakable, width=0.92\linewidth,
                  colback=gray!3, colframe=gray!45, boxrule=0.5pt,
                  arc=3pt, left=12pt, right=12pt, top=8pt, bottom=8pt,
                  title={Argument flow \quad
                         \disc{v}\,V \;\disc{p}\,P \;\disc{h}\,H \;\disc{c}\,C}]

\sknode{\disc{v}}{Assumption~\ref{ass:cost}}{entry costs a little, and ties split evenly}{%
  $0<\kappa<1/n$ per front opened --- worth opening one, never worth opening all\sktag{setup}}

\skstep{so}

\sknode{\disc{p}}{Proposition~\ref{prop:exist}}{the contest has a mixed-strategy equilibrium}{%
  for any opened fronts, players randomize how hard to fight on each\sktag{engine}}

\skstep{and}

\sknode{\disc{p}}{Lemma~\ref{lem:concentrate}}{each side opens only a few fronts}{%
  at most $\lceil 1/(2\kappa)\rceil$ of them --- a bounded number, not all $n$\sktag{engine}}

\skstep{so}

\sknode{\disc{h}}{Corollary~\ref{cor:sparse}}{those few fronts don't grow with $n$}{%
  the support stays size $O(1/\kappa)$, independent of $n$\sktag{hinge}}

\skstep{so}

\sknode{\disc{c}}{Theorem~\ref{thm:char}}{exact closed-form for how they play}{%
  a symmetric mixed equilibrium: open $m^\star=\lceil 1/(2\kappa)\rceil$ fronts,
  spend uniformly on $[0,2/m^\star]$ each
  \;\textnormal{(needs Proposition~\ref{prop:exist})}\sktag{main result}}

\skstep{still open}

\sknode{\disc{c}}{Uniqueness?}{is this the only solution?}{%
  rule out asymmetric equilibria on the same support --- would turn
  Theorem~\ref{thm:char} green\sktag{open frontier}}

\medskip
\hrule
\smallskip
{\footnotesize\itshape
Read top-to-bottom: each step follows from the ones above; the word on each
$\downarrow$ is the logic. ``(needs $X$)'' flags a step that leans on something
not directly above it. Disc colour = current verification status.}

\end{tcolorbox}
\end{center}
 from prototype.tex (after the appendix).
% =====================================================================
\clearpage
\section*{Dependency-flow skeleton}
\addcontentsline{toc}{section}{Dependency-flow skeleton}

% A NODE brick: disc, full name (bold), plain-English role; then an indented
% smaller concrete-detail line. #1 disc | #2 \ref name | #3 role | #4 detail.
\newcommand{\sknode}[4]{%
  \noindent #1\, \textbf{#2} \;---\; #3\par
  \nobreak\smallskip
  {\footnotesize\color{gray!75}\hspace*{1.6em}#4\par}}

% The logic arrow between steps, carrying the connective word.
\newcommand{\skstep}[1]{%
  \par\medskip\centerline{$\downarrow$\ \footnotesize\itshape\color{gray!70}#1}\par\medskip}

% A faint right-aligned tier tag (optional, must not dominate).
\newcommand{\sktag}[1]{\hfill{\scriptsize\color{gray!45}#1}}

\begin{center}
\begin{tcolorbox}[enhanced, breakable, width=0.92\linewidth,
                  colback=gray!3, colframe=gray!45, boxrule=0.5pt,
                  arc=3pt, left=12pt, right=12pt, top=8pt, bottom=8pt,
                  title={Argument flow \quad
                         \disc{v}\,V \;\disc{p}\,P \;\disc{h}\,H \;\disc{c}\,C}]

\sknode{\disc{v}}{Assumption~\ref{ass:cost}}{entry costs a little, and ties split evenly}{%
  $0<\kappa<1/n$ per front opened --- worth opening one, never worth opening all\sktag{setup}}

\skstep{so}

\sknode{\disc{p}}{Proposition~\ref{prop:exist}}{the contest has a mixed-strategy equilibrium}{%
  for any opened fronts, players randomize how hard to fight on each\sktag{engine}}

\skstep{and}

\sknode{\disc{p}}{Lemma~\ref{lem:concentrate}}{each side opens only a few fronts}{%
  at most $\lceil 1/(2\kappa)\rceil$ of them --- a bounded number, not all $n$\sktag{engine}}

\skstep{so}

\sknode{\disc{h}}{Corollary~\ref{cor:sparse}}{those few fronts don't grow with $n$}{%
  the support stays size $O(1/\kappa)$, independent of $n$\sktag{hinge}}

\skstep{so}

\sknode{\disc{c}}{Theorem~\ref{thm:char}}{exact closed-form for how they play}{%
  a symmetric mixed equilibrium: open $m^\star=\lceil 1/(2\kappa)\rceil$ fronts,
  spend uniformly on $[0,2/m^\star]$ each
  \;\textnormal{(needs Proposition~\ref{prop:exist})}\sktag{main result}}

\skstep{still open}

\sknode{\disc{c}}{Uniqueness?}{is this the only solution?}{%
  rule out asymmetric equilibria on the same support --- would turn
  Theorem~\ref{thm:char} green\sktag{open frontier}}

\medskip
\hrule
\smallskip
{\footnotesize\itshape
Read top-to-bottom: each step follows from the ones above; the word on each
$\downarrow$ is the logic. ``(needs $X$)'' flags a step that leans on something
not directly above it. Disc colour = current verification status.}

\end{tcolorbox}
\end{center}
 from prototype.tex (after the appendix).
% =====================================================================
\clearpage
\section*{Dependency-flow skeleton}
\addcontentsline{toc}{section}{Dependency-flow skeleton}

% A NODE brick: disc, full name (bold), plain-English role; then an indented
% smaller concrete-detail line. #1 disc | #2 \ref name | #3 role | #4 detail.
\newcommand{\sknode}[4]{%
  \noindent #1\, \textbf{#2} \;---\; #3\par
  \nobreak\smallskip
  {\footnotesize\color{gray!75}\hspace*{1.6em}#4\par}}

% The logic arrow between steps, carrying the connective word.
\newcommand{\skstep}[1]{%
  \par\medskip\centerline{$\downarrow$\ \footnotesize\itshape\color{gray!70}#1}\par\medskip}

% A faint right-aligned tier tag (optional, must not dominate).
\newcommand{\sktag}[1]{\hfill{\scriptsize\color{gray!45}#1}}

\begin{center}
\begin{tcolorbox}[enhanced, breakable, width=0.92\linewidth,
                  colback=gray!3, colframe=gray!45, boxrule=0.5pt,
                  arc=3pt, left=12pt, right=12pt, top=8pt, bottom=8pt,
                  title={Argument flow \quad
                         \disc{v}\,V \;\disc{p}\,P \;\disc{h}\,H \;\disc{c}\,C}]

\sknode{\disc{v}}{Assumption~\ref{ass:cost}}{entry costs a little, and ties split evenly}{%
  $0<\kappa<1/n$ per front opened --- worth opening one, never worth opening all\sktag{setup}}

\skstep{so}

\sknode{\disc{p}}{Proposition~\ref{prop:exist}}{the contest has a mixed-strategy equilibrium}{%
  for any opened fronts, players randomize how hard to fight on each\sktag{engine}}

\skstep{and}

\sknode{\disc{p}}{Lemma~\ref{lem:concentrate}}{each side opens only a few fronts}{%
  at most $\lceil 1/(2\kappa)\rceil$ of them --- a bounded number, not all $n$\sktag{engine}}

\skstep{so}

\sknode{\disc{h}}{Corollary~\ref{cor:sparse}}{those few fronts don't grow with $n$}{%
  the support stays size $O(1/\kappa)$, independent of $n$\sktag{hinge}}

\skstep{so}

\sknode{\disc{c}}{Theorem~\ref{thm:char}}{exact closed-form for how they play}{%
  a symmetric mixed equilibrium: open $m^\star=\lceil 1/(2\kappa)\rceil$ fronts,
  spend uniformly on $[0,2/m^\star]$ each
  \;\textnormal{(needs Proposition~\ref{prop:exist})}\sktag{main result}}

\skstep{still open}

\sknode{\disc{c}}{Uniqueness?}{is this the only solution?}{%
  rule out asymmetric equilibria on the same support --- would turn
  Theorem~\ref{thm:char} green\sktag{open frontier}}

\medskip
\hrule
\smallskip
{\footnotesize\itshape
Read top-to-bottom: each step follows from the ones above; the word on each
$\downarrow$ is the logic. ``(needs $X$)'' flags a step that leans on something
not directly above it. Disc colour = current verification status.}

\end{tcolorbox}
\end{center}
